\scanpage{8}
de $\underline O_X$-Module provenant de l'extrême droite
[NB. les structures extrême-droites des $Q^i[M]$ sont compatibles avec les op. simpliciales, i.e. $Q^{\unsure{i}}[M]$ est un Module cosimplicial à $\underline O_X$ opérant par ses op. extrême-droites ; mais bien sûr cette dernière structure n'est pas conservée par les hom $M-N$ qui sont non $\underline O_X$-lin (mais seult des op. diff. rel / $S$)]. De $\overline{\mathrm m}$, un hom $D:M-N$ de $\operatorname{Mod}(X/S)$ est récupéré à partir de $Q^0[D]:Q^0[M]\unsure{-|\longrightarrow}Q^0[N]$ comme l'hom induit sur les ss-faisceaux des \deleted{\unsure{objets}}\inserted{sections} stationnaires.

4. Une façon équivalente d'obtenir $Q^*[M]$ est de considérer
\diagram{\begin{tikzcd}[ampersand replacement=\&,column sep=large]
X^{\widehat{[2]}}\arrow[r,shift left=.7ex]\arrow[r,shift right=.7ex]\&X^{\widehat{[1]}}\arrow[r,"\widehat{\operatorname{pr}}_2"]\&X^{\widehat{[0]}}=X
\end{tikzcd}}
et les produits fibrés [dans la cat. des présch. formels sur $X$], \deleted{\unsure{donc}} c'est aussi
\margin{7}
